\Needspace{12\baselineskip}
\begin{center}\textit{Ein allgemeiner (fundamentaler) Satz wird anhand von Zahlenbeispielen aufgestellt; daraus werden Folgerungen gezogen.}\end{center}

\Article{130}
Nachdem wir streng bewiesen haben, dass jede Primzahl der Form $4n+1$, sowohl positiv als auch negativ genommen, ein Nichtrest einer kleineren Primzahl ist, \SourceLeaf{104}gehen wir nun zu einem genaueren und allgemeineren Vergleich der Primzahlen über, je nachdem, ob eine Rest oder Nichtrest einer anderen ist.

Oben haben wir streng bewiesen, dass $-3$ und $+5$ Reste oder Nichtreste aller Primzahlen sind, die jeweils Reste oder Nichtreste von $3, 5$ sind.

Die Untersuchung der folgenden Zahlen zeigt, dass $-7$, $-11$, $+13$, $+17$, $-19$, $-23$, $+29$, $-31$, $+37$, $+41$, $-43$, $-47$, $+53$, $-59$ usw. Reste oder Nichtreste aller Primzahlen sind, die, positiv genommen, jeweils Reste oder Nichtreste jener Primzahlen sind. Diese Untersuchung lässt sich mit Tabelle II sehr leicht durchführen.

Bei etwas Aufmerksamkeit wird jeder bemerken, dass unter diesen Primzahlen diejenigen der Form $4n+1$ ein positives Vorzeichen haben, diejenigen der Form $4n+3$ dagegen ein negatives.

\Article{131}
Wir werden bald beweisen, dass das, was wir hier anhand von Beispielen gefunden haben, allgemein gilt. Bevor wir diese Aufgabe angehen, müssen wir jedoch alles herleiten, was aus dem Satz folgt, wenn man ihn als wahr voraussetzt. Den Satz selbst sprechen wir so aus.

\emph{Ist $p$ eine Primzahl der Form $4n+1$, so ist $+p$, und hat $p$ die Form $4n+3$, so ist $-p$ ein Rest oder Nichtrest jeder Primzahl, die, positiv genommen, ein Rest oder Nichtrest von $p$ ist.}

Da fast alles, was sich über quadratische Reste sagen lässt, auf diesem Satz beruht, ist die Bezeichnung \emph{Fundamentalsatz}, die wir von nun an verwenden werden, durchaus passend.

Um unsere Schlüsse möglichst kurz darzustellen, bezeichnen wir mit $a, a', a''$ usw. Primzahlen der Form $4n+1$, mit $b, b', b''$ usw. Primzahlen der Form $4n+3$; mit $A, A', A''$ usw. beliebige Zahlen der Form $4n+1$, mit $B, B', B''$ usw. dagegen beliebige Zahlen der Form $4n+3$. Schließlich soll der Buchstabe $R$ zwischen zwei Größen bedeuten, dass die erste ein Rest der zweiten ist; der Buchstabe $N$ soll die entgegengesetzte Bedeutung haben. \emph{Zum Beispiel} bedeutet $+5R11$, $\pm 2N5$, dass $+5$ ein Rest von $11$ und $+2$ oder $-2$ ein Nichtrest von $5$ ist. \SourceLeaf{105}Verbindet man nun den Fundamentalsatz mit den Sätzen aus Artikel $111$, so erhält man leicht die folgenden Aussagen.

\[
\begin{array}{rccc}&\text{Wenn}&&\text{dann}\\[5pt]
1.&\pm aRa'&\mathinner{.\ .\ .\ .\ .\ .\ .}&\pm a' Ra\\[4pt]
2.&\pm aNa'&\mathinner{.\ .\ .\ .\ .\ .\ .}&\pm a' Na\\[4pt]
3.&\left\{\begin{array}{l}+aRb\\-aNb\end{array}\right\}&\mathinner{.\ .\ .\ .\ .\ .\ .}&\pm bRa\\[4pt]
4.&\left\{\begin{array}{l}+aNb\\-aRb\end{array}\right\}&\mathinner{.\ .\ .\ .\ .\ .\ .}&\pm bNa\\[4pt]
5.&\pm bRa&\mathinner{.\ .\ .\ .\ .\ .\ .}&\left\{\begin{array}{l}+aRb\\-aNb\end{array}\right.\\[4pt]
6.&\pm bNa&\mathinner{.\ .\ .\ .\ .\ .\ .}&\left\{\begin{array}{l}+aNb\\-aRb\end{array}\right.\\[4pt]
7.&\left\{\begin{array}{l}+bRb'\\-bNb'\end{array}\right\}&\mathinner{.\ .\ .\ .\ .\ .\ .}&\left\{\begin{array}{l}+b' Nb\\-b' Rb\end{array}\right.\\[4pt]
8.&\left\{\begin{array}{l}+bNb'\\-bRb'\end{array}\right\}&\mathinner{.\ .\ .\ .\ .\ .\ .}&\left\{\begin{array}{l}+b' Rb\\-b' Nb\end{array}\right.\\[4pt]
\end{array}
\]

\Article{132}
Diese enthalten alle Fälle, die beim Vergleich zweier Primzahlen auftreten können. Die folgenden Aussagen beziehen sich auf beliebige Zahlen; ihre Beweise liegen jedoch weniger auf der Hand.

\[
\begin{array}{rccc}&\text{Wenn}&&\text{dann}\\[5pt]
9.&\pm aRA&\mathinner{.\ .\ .\ .\ .\ .\ .}&\pm ARa\\[4pt]
10.&\pm bRA&\mathinner{.\ .\ .\ .\ .\ .\ .}&\left\{\begin{array}{l}+ARb\\-ANb\end{array}\right.\\[4pt]
11.&+aRB&\mathinner{.\ .\ .\ .\ .\ .\ .}&\pm BRa\\[4pt]
12.&-aRB&\mathinner{.\ .\ .\ .\ .\ .\ .}&\pm BNa\\[4pt]
13.&+bRB&\mathinner{.\ .\ .\ .\ .\ .\ .}&\left\{\begin{array}{l}-BRb\\+BNb\end{array}\right.\\[4pt]
14.&-bRB&\mathinner{.\ .\ .\ \phantom{.\ .\ .}\ .}&\left\{\begin{array}{l}+BRb\\-BNb\end{array}\right.\\[4pt]
\end{array}
\]

\SourceLeaf{106}Da die Beweise all dieser Aussagen aus denselben Prinzipien folgen, müssen wir sie nicht alle ausführen. Der Beweis von Aussage $9$, den wir angeben, kann als Beispiel dienen. Zunächst ist zu beachten, dass jede Zahl der Form $4n+1$ entweder keinen Faktor der Form $4n+3$ hat oder zwei oder vier usw.; \emph{das heißt,} die Anzahl solcher Faktoren (unter denen auch gleiche vorkommen können) ist stets gerade. Jede Zahl der Form $4n+3$ enthält dagegen eine ungerade Anzahl von Faktoren der Form $4n+3$ (\emph{das heißt} entweder einen oder drei oder fünf usw.). Die Anzahl der Faktoren der Form $4n+1$ bleibt unbestimmt.

\emph{Aussage} $9$ wird so bewiesen. Sei $A$ das Produkt der Primfaktoren $a',a'',a'''$ usw., $b,b',b''$ usw. Dann ist die Anzahl der Faktoren $b,b',b''$ usw. gerade (es können auch keine vorkommen, was auf dasselbe hinausläuft). Ist nun $a$ ein Rest von $A$, so auch von allen Faktoren $a',a'',a'''$ usw., $b,b',b''$ usw. Nach den Aussagen $1,3$ des vorigen Artikels ist daher jeder dieser Faktoren ein Rest von $a$, also auch ihr Produkt $A$. Dasselbe muss für $-A$ gelten. --- Ist dagegen $-a$ ein Rest von $A$ und damit auch von allen Faktoren $a',a''$ usw., $b,b'$ usw., so sind die einzelnen $a',a''$ usw. Reste von $a$, die einzelnen $b,b'$ usw. dagegen Nichtreste. Da die Anzahl der letzteren gerade ist, ist das Produkt aller Faktoren, \emph{das heißt} $A$, ein Rest von $a$ und somit auch $-A$.

\Article{133}
Stellen wir die Untersuchung noch allgemeiner an. Betrachten wir zwei beliebige teilerfremde ungerade Zahlen mit beliebigen Vorzeichen, $P$ und $Q$. Man denke sich $P$ ohne Rücksicht auf sein Vorzeichen in seine Primfaktoren zerlegt und bezeichne mit $p$, wie viele davon $Q$ als Nichtrest haben. Kommt eine Primzahl, von der $Q$ ein Nichtrest ist, mehrmals unter den Faktoren von $P$ vor, so muss sie auch mehrmals gezählt werden. Entsprechend sei $q$ die Anzahl der Primfaktoren von $Q$, von denen $P$ ein Nichtrest ist. Dann stehen die Zahlen $p,q$ in einer bestimmten Beziehung zueinander, die von der Art der Zahlen $P,Q$ abhängt. Ist nämlich eine der Zahlen $p,q$ gerade oder ungerade, so zeigt die Form der Zahlen $P,Q$, ob die andere gerade oder ungerade ist. Die folgende Tabelle stellt diese Beziehung dar.

$p,q$ sind beide gerade oder beide ungerade, wenn die Zahlen $P,Q$ die Formen haben:

\[
\begin{array}{rl}1.&+A,\quad+A'\\2.&+A,\quad-A'\\\SourceLeaf{107}3.&+A,\quad+B\\4.&+A,\quad-B\\5.&-A,\quad-A'\\6.&+B,\quad-B'\end{array}
\]

Dagegen ist von den Zahlen $p,q$ eine gerade und die andere ungerade, wenn die Zahlen $P,Q$ die Formen haben:

\[
\begin{array}{rl}7.&-A,\quad+B\\8.&-A,\quad-B\\9.&+B,\quad+B'\\10.&-B,\quad-B'{}^{*)}\end{array}
\]

\HistoricalDisplayNote{*)}{\begin{tabular}{@{}ll@{}}Sei &$l=1$, wenn beide $P,Q\equiv3\,(\operatorname{mod.}4)$, sonst $l=0$\\&$m=1$, wenn beide $P,Q$ negativ sind, sonst $m=0$\end{tabular}\par Dann hängt die Beziehung von $l+m$ ab.}

\emph{Beispiel.} Die gegebenen Zahlen seien $-55$ und $+1197$, die zum vierten Fall gehören. Dabei ist $1197$ ein Nichtrest eines Primfaktors von $55$. Dieser Faktor ist die Zahl $5$, während $-55$ ein Nichtrest dreier Primfaktoren von $1197$ ist, nämlich der Zahlen $3,3,19$.

Bezeichnen $P$ und $Q$ Primzahlen, so gehen diese Aussagen in die aus Artikel $131$ über. Hier können $p$ und $q$ nicht größer als $1$ sein. Wird also $p$ als gerade angenommen, so muss es $=0$ sein, \emph{das heißt,} $Q$ ist ein Rest von $P$; ist dagegen $p$ ungerade, so ist $Q$ ein Nichtrest von $P$. Und umgekehrt. Schreibt man also $a,b$ anstelle von $A,B$, so folgt aus $8$: Ist $-a$ ein Rest oder Nichtrest von $b$, so ist $-b$ ein Nichtrest oder Rest von $a$. Das stimmt mit $3$ und $4$ aus Artikel $131$ überein.

Allgemein kann $Q$ aber nur dann ein Rest von $P$ sein, wenn $p=0$. Ist also $p$ ungerade, so ist $Q$ sicher ein Nichtrest von $P$.

Hieraus lassen sich auch die Aussagen des vorigen Artikels ohne Schwierigkeit herleiten.

Bald wird sich außerdem zeigen, dass diese allgemeine Darstellung mehr als eine unfruchtbare Spekulation ist, da sich ein vollständiger Beweis des Fundamentalsatzes ohne sie kaum führen lässt.

\SourceLeaf{108}\Article{134}
Wir gehen nun zur Herleitung dieser Aussagen über.

I.\quad Man denke sich wie zuvor $P$ ohne Berücksichtigung der Vorzeichen in seine Primfaktoren zerlegt. Auch $Q$ werde auf beliebige Weise in Faktoren zerlegt, wobei aber das Vorzeichen von $Q$ berücksichtigt wird. Man bilde Paare aus jedem der ersten Faktoren und jedem der zweiten. Bezeichnet $s$ die Anzahl aller Paare, in denen der Faktor von $Q$ ein Nichtrest des Faktors von $P$ ist, so sind $p$ und $s$ beide gerade oder beide ungerade. Denn die Primfaktoren von $P$ seien $f,f',f''$ usw., und unter den Faktoren, in die $Q$ zerlegt ist, seien $m$ Nichtreste von $f$, $m'$ Nichtreste von $f'$, $m''$ Nichtreste von $f''$ usw. Dann sieht man leicht, dass

\[
s=m+m'+m''+\text{usw.}
\]

gilt, während $p$ angibt, wie viele der Zahlen $m,m',m''$ usw. ungerade sind. Daraus folgt unmittelbar, dass $s$ gerade ist, wenn $p$ gerade ist, und ungerade, wenn $p$ ungerade ist.

II.\quad Dies gilt allgemein, auf welche Weise $Q$ auch in Faktoren zerlegt wurde. Gehen wir zu besonderen Fällen über. Betrachten wir zunächst die Fälle, in denen die eine Zahl, $P$, positiv ist, die andere, $Q$, dagegen entweder die Form $+A$ oder die Form $-B$ hat. Man zerlege $P,Q$ in ihre Primfaktoren, gebe jedem Faktor von $P$ ein positives Vorzeichen und jedem Faktor von $Q$ ein positives oder negatives, je nachdem, ob er die Form $a$ oder $b$ hat. Dann hat $Q$ offensichtlich entweder die Form $+A$ oder $-B$, wie verlangt. Man bilde Paare aus jedem Faktor von $P$ und jedem Faktor von $Q$, und bezeichne wie zuvor mit $s$ die Anzahl der Paare, in denen der Faktor von $Q$ ein Nichtrest des Faktors von $P$ ist. Entsprechend sei $t$ die Anzahl der Paare, in denen der Faktor von $P$ ein Nichtrest des Faktors von $Q$ ist. Aus dem Fundamentalsatz folgt, dass diese beiden Mengen von Paaren übereinstimmen, also $s=t$. Aus dem soeben Bewiesenen folgt schließlich $p\equiv s\, (\mathrm{mod.}\ 2)$, $q\equiv t\, (\mathrm{mod.}\ 2)$ und damit $p\equiv q\, (\mathrm{mod.}\ 2)$.

Wir haben damit die Aussagen $1,3,4$ und $6$ aus Artikel $133$.

Die übrigen Aussagen lassen sich auf ähnliche Weise direkt gewinnen, erfordern aber eine neue Überlegung. Leichter lassen sie sich aus dem Vorhergehenden so herleiten.

\Needspace{12\baselineskip}
III.\quad Wieder seien $P,Q$ beliebige teilerfremde ungerade Zahlen, und $p,q$ die Anzahlen der Primfaktoren von $P,Q$, von denen $Q$ bzw. $P$ Nichtreste sind. Schließlich sei $p'$ die Anzahl der Primfaktoren von $P$, von denen \SourceLeaf{109}$-Q$ ein Nichtrest ist (ist $Q$ selbst negativ, so bezeichnet $-Q$ offensichtlich eine positive Zahl). Man teile nun alle Primfaktoren von $P$ in vier Klassen ein:

1) Faktoren der Form $a$, von denen $Q$ ein Rest ist.

2) Faktoren der Form $b$, von denen $Q$ ein Rest ist. Ihre Anzahl sei $\chi$.

3) Faktoren der Form $a$, von denen $Q$ ein Nichtrest ist. Ihre Anzahl sei $\psi$.

4) Faktoren der Form $b$, von denen $Q$ ein Nichtrest ist. Ihre Anzahl $=\omega$.

Dann sieht man leicht, dass $p=\psi+\omega$, $p'=\chi+\psi$ gilt.

Hat nun $P$ die Form $\pm A$, so ist $\chi+\omega$ und damit auch $\chi-\omega$ gerade. Also gilt $p'=p+\chi-\omega\equiv p\, (\mathrm{mod.}\ 2)$. Hat dagegen $P$ die Form $\pm B$, so zeigt ein ähnlicher Schluss, dass die Zahlen $p,p'$ nach dem Modul $2$ inkongruent sind.

IV.\quad Wenden wir dies auf die einzelnen Fälle an. Zunächst hätten sowohl $P$ als auch $Q$ die Form $+A$. Dann gilt nach Aussage $1$ $p\equiv q\, (\mathrm{mod.}\ 2)$; es gilt aber auch $p'\equiv p\, (\mathrm{mod.}\ 2)$, also $p'\equiv q\ (\mathrm{mod.}\ 2)$. Das stimmt mit Aussage $2$ überein. --- Entsprechend habe $P$ die Form $-A$. Hat $Q$ die Form $+A$, so gilt $p\equiv q\, (\mathrm{mod.}\ 2)$ nach der soeben bewiesenen Aussage $2$. Da $p'\equiv p$, folgt daraus $p'\equiv q$. Damit ist auch Aussage $5$ bewiesen.

Auf dieselbe Weise folgt Aussage $7$ aus $3$; Aussage $8$ aus $4$ oder $7$; Aussage $9$ aus $6$ und aus derselben Aussage auch $10$.
